% Generated by Listas/make_statements.py from pratica-lecture-01.md, item 14. Do not edit.
\textbf{Source.} Muñoz-Garcia \& Toro-Gonzalez (2019), \emph{Strategy and Game Theory: Practice Exercises with Answers}, Ch. 1, Exercise 12 ``IDSDS and Rationality'', p.~23 {[}PDF 36{]}. Reproduced verbatim (statement only; the book's worked answer follows it).



\begin{quote}

Consider the following \(3 \times 3\) matrix (Fig. 1.30):



{\def\LTcaptype{none} % do not increment counter

\begin{longtable}[]{@{}llll@{}}

\toprule\noalign{}

Player 1 ~Player 2 & \(F\) & \(C\) & \(B\) \\

\midrule\noalign{}

\endhead

\bottomrule\noalign{}

\endlastfoot

\(F\) & 0, 5 & 2, 4 & 2, 3 \\

\(C\) & 2, 3 & 0, 5 & 3, 2 \\

\(B\) & 5, 0 & 2, 4 & 2, 3 \\

\end{longtable}

}



\emph{Fig. 1.30} Normal-form representation of a three-players game



\begin{enumerate}

\def\labelenumi{\alph{enumi}.}

\item

  Does any player have a strictly dominant strategy?

\item

  Does any player have a strictly dominated strategy?

\item

  Use your results from part (b) to delete players' strictly dominated strategies assuming rationality, but not assuming common knowledge of rationality. {[}Hint: This means that every player can delete the strategies he finds to be strictly dominated but cannot put himself in the shoes of his opponent to delete the strategies his opponent finds to be strictly dominated.{]}

\item

  Use your results from part (b) to delete players' strictly dominated strategies assuming common knowledge of rationality. {[}Hint: This means that players can put themselves in the shoes of their opponents to delete strategies each other finds to be strictly dominated, as in the standard approach of IDSDS.{]}

\end{enumerate}

\end{quote}

